\section{Physical output after exchanging auxiliary registers}
\label{sec:peps_physical_first_output}

The source replacement of
\cite[proof of Theorem~5.2]{OpenAI2026PolynomialPEPS},
\texttt{04-compression.tex}, lines 199--227, has output memory
$\mathcal A\otimes(\mathcal H\otimes\mathcal E)$, where $\mathcal A$
contains the retained auxiliary registers, $\mathcal H$ contains the
original physical registers, and $\mathcal E$ contains the original
discarded private registers. Write $\tau$ for the canonical isometry
from this memory to $\mathcal H\otimes(\mathcal A\otimes\mathcal E)$.

\begin{definition}[Physical-first replacement]
    \label{def:peps_physical_first_replacement}
    \lean{TNLean.PEPS.PairEffect.SourceCircuit.physicalFirstOutput}
    \lean{TNLean.PEPS.PairEffect.EffectCircuit.physicalFirstReplacement}
    \leanok
    If $R$ is the constructed source-only replacement, define
    $R_{\rm phys}=\tau R$. The exchange acts on pure tensors by
    \begin{align}
        \tau\bigl(u\otimes(p\otimes v)\bigr)
        =p\otimes(u\otimes v).
        \label{eq:peps_physical_first_exchange}
    \end{align}
\end{definition}

\begin{theorem}[Preservation of gate and source occurrences]
    \label{thm:peps_physical_first_occurrences}
    \lean{TNLean.PEPS.PairEffect.SourceCircuit.sourceOrder_physicalFirstOutput}
    \lean{TNLean.PEPS.PairEffect.SourceCircuit.physicalFirstLocationsEquiv}
    \lean{TNLean.PEPS.PairEffect.SourceCircuit.physicalFirstSourcesEquiv}
    \lean{TNLean.PEPS.PairEffect.SourceCircuit.participants_physicalFirstLocationsEquiv}
    \lean{TNLean.PEPS.PairEffect.SourceCircuit.branchCoefficient_physicalFirstLocationsEquiv}
    \lean{TNLean.PEPS.PairEffect.SourceCircuit.endpoints_physicalFirstSourcesEquiv}
    \lean{TNLean.PEPS.PairEffect.SourceCircuit.sourceVectorAt_physicalFirstSourcesEquiv}
    \lean{TNLean.PEPS.PairEffect.SourceCircuit.participationCount_physicalFirstOutput}
    \lean{TNLean.PEPS.PairEffect.SourceCircuit.isExpansionBounded_physicalFirstOutput}
    \leanok
    \uses{def:peps_physical_first_replacement}
    Appending $\tau$ preserves the gate occurrences, source occurrences,
    participating sets, branch-label types, branch coefficients, endpoint spaces and source
    vectors of $R$. It preserves source order, each party's total number
    of participating occurrences, and the bounds on monomial counts and
    coefficient sums.
\end{theorem}

\begin{proof}
    \lean{TNLean.PEPS.PairEffect.SourceCircuit.branchLabels_physicalFirstLocationsEquiv}
    \lean{TNLean.PEPS.PairEffect.SourceCircuit.card_gateLocations_physicalFirstOutput}
    \lean{TNLean.PEPS.PairEffect.SourceCircuit.card_sourceLocations_physicalFirstOutput}
    \lean{TNLean.PEPS.PairEffect.SourceCircuit.eval_exchangeBlocks_operator}
    \lean{TNLean.PEPS.PairEffect.SourceCircuit.eval_exchangeBlocks_threeFactors}
    \lean{TNLean.PEPS.PairEffect.SourceCircuit.eval_exchangeBlocks_triple_tmul}
    \lean{TNLean.PEPS.PairEffect.SourceCircuit.isAllowed_physicalFirstOutput}
    \lean{TNLean.PEPS.PairEffect.SourceCircuit.physicalFirstParticipationEquiv}
    \lean{TNLean.PEPS.PairEffect.SourceCircuit.slotCount_physicalFirstLocationsEquiv}
    \lean{TNLean.PEPS.PairEffect.SourceCircuit.sourceDims_physicalFirstSourcesEquiv}
    \lean{TNLean.PEPS.PairEffect.SourceCircuit.sourceOrder_exchangeBlocks}
    \leanok
    The exchange consists only of permutations of register order and
    introduces no gate or source occurrence. Include each occurrence of
    $R$ into the first part of the composition. The defining data at
    these occurrences are unchanged, and the second part contributes
    empty occurrence sets. The gate and source cardinalities are therefore unchanged,
    and allowedness is equivalent before and after the exchange. Its three-factor
    tensor identity extends by linearity to the full operator identity.
\end{proof}

\begin{theorem}[Finite-dimensional discarded registers]
    \label{thm:peps_auxiliary_register_finiteness}
    \lean{TNLean.PEPS.PairEffect.EffectCircuit.finiteDimensional_auxiliary_registers}
    \lean{TNLean.PEPS.PairEffect.EffectCircuit.finiteDimensional_auxiliary_private_registers}
    \lean{TNLean.PEPS.PairEffect.EffectCircuit.finiteDimensional_auxiliary_private_mem}
    \leanok
    Every auxiliary register constructed by the source replacement is
    finite-dimensional. If every original discarded private register is
    finite-dimensional, the same holds for every register of
    $\mathcal A\otimes\mathcal E$ and for their tensor product. No
    numerical bound on private dimensions is required.
\end{theorem}

\begin{proof}
    \lean{TNLean.PEPS.PairEffect.Layout.finiteDimensional_registers_append}
    \lean{TNLean.PEPS.PairEffect.Layout.finiteDimensional_registers_mapOwner}
    \lean{TNLean.PEPS.PairEffect.PartyChain.finiteDimensional_registers_sources_prepStack}
    \lean{TNLean.PEPS.PairEffect.finiteDimensional_registers_sources_comp}
    \lean{TNLean.PEPS.PairEffect.finiteDimensional_registers_sources_prepGate}
    \lean{TNLean.PEPS.PairEffect.finiteDimensional_registers_sources_prepWord}
    \leanok
    A prepared pair stack consists of finitely many finite-dimensional
    Euclidean registers. Induction on the stack length, on the monomial,
    on the gate expansion, and on the original chronology proves the
    assertion for each auxiliary register. Changing an owner leaves the
    register space unchanged, and concatenation preserves this assertion.
    Finally, a finite tensor product of finite-dimensional spaces is
    finite-dimensional.
\end{proof}

\begin{theorem}[Physical partial trace after the exchange]
    \label{thm:peps_physical_first_trace}
    \lean{TNLean.PEPS.PairEffect.physicalFirstReadout_exchange_apply}
    \lean{TNLean.PEPS.PairEffect.physicalFirstReadout_exchange_outerProduct}
    \lean{TNLean.PEPS.PairEffect.partialTrace_physicalFirstReadout_exchange}
    \lean{TNLean.PEPS.PairEffect.SourceCircuit.physicalDensity_physicalFirstOutput}
    \leanok
    \uses{def:peps_physical_first_replacement}
    Choose orthonormal coordinates indexed by $U$, $X$ and $V$ on
    $\mathcal A$, $\mathcal H$ and $\mathcal E$, respectively. For
    $z\in\mathcal A\otimes(\mathcal H\otimes\mathcal E)$, let
    $z_{i,(k,j)}$ denote its coordinates when the physical coordinate
    precedes the private and auxiliary coordinates. The coordinates of
    $\tau z$ are $(\tau z)_{i,(j,k)}=z_{i,(k,j)}$. Consequently,
    \begin{align}
        \operatorname{Tr}_{\mathcal A\otimes\mathcal E}
        |\tau z\rangle\langle\tau z|
        =\operatorname{Tr}_{\mathcal E\otimes\mathcal A}
        |z\rangle\langle z|.
        \label{eq:peps_physical_first_trace}
    \end{align}
    This identity holds for every vector, without a normalization premise.
\end{theorem}

\begin{proof}
    \lean{TNLean.PEPS.PairEffect.SourceCircuit.physicalDensity_pure_input}
    \lean{TNLean.PEPS.PairEffect.exchangedGarbageBasis}
    \lean{TNLean.PEPS.PairEffect.extendedPhysicalReadout_triple_tmul}
    \lean{TNLean.PEPS.PairEffect.physicalFirstReadout}
    \lean{TNLean.PEPS.PairEffect.physicalFirstReadout_tmul}
    \leanok
    Formula~(\ref{eq:peps_physical_first_exchange}) gives the coordinate
    identity on pure tensors. Tensor-product induction extends it to
    arbitrary vectors. Their outer-product matrices are related by the
    simultaneous reindexing $(i,(j,k))\mapsto(i,(k,j))$. Reindexing the
    traced coordinate by the bijection $U\times V\to V\times U$
    preserves the partial trace. The physical-first coordinate isometry uses the
    tensor product of the physical basis and the auxiliary--private basis.
    Evaluation of a pure input has density equal to the outer product of its
    evaluated vector in these coordinates. Applying this to the output vector of
    the composed circuit proves the final assertion.
\end{proof}

\begin{theorem}[Exact sources recover the replacement physical state]
    \label{thm:peps_physical_first_exact_sources}
    \lean{TNLean.PEPS.PairEffect.EffectCircuit.physicalDensity_physicalFirstReplacement}
    \lean{TNLean.PEPS.PairEffect.EffectCircuit.physicalSourceReplacedDensity_physicalFirstReplacement_exact}
    \lean{TNLean.PEPS.PairEffect.OriginalCircuit.rectangularTraceNorm_physicalFirst_density_sub_le_half}
    \leanok
    \uses{thm:peps_physical_first_trace,thm:peps_source_only_density_error}
    Let $W$ be an original circuit with at most $M$ nonprivate gates,
    satisfying the original expansion bounds. Set
    $\delta=\varepsilon/(8\max\{1,M\})$ for $\varepsilon>0$, and
    construct its replacement $R$ using this budget. For a pure input
    $\psi$ with $\|\psi\|\le1$, let $\rho_W$ be the original physical
    output state. In the source-replacement expression for $R_{\rm phys}$,
    insert the exact operator of each original pair source. The resulting
    physical operator is the replacement state $\rho_R$, and
    $\|\rho_R-\rho_W\|_1\le\varepsilon/2$.
    The physical coordinates may be the labelled tensor coordinates with
    an arbitrary original local dimension $d_p$ at each party $p$.
\end{theorem}

\begin{proof}\leanok
    Exact source operators recover the output operator of the source-only
    circuit. Formula~(\ref{eq:peps_physical_first_trace}) identifies its
    physical partial trace with $\rho_R$. Apply
    Theorem~\ref{thm:peps_source_only_density_error} at the original
    gate-count bound $M$.
\end{proof}
